\section{AMC 10 10A 2007}
        \begin{enumerate}
        	\item (\textit{AMC 10 10A 2007}) One ticket to a show costs $\$20$ at full price. Susan buys 4 tickets using a coupon that gives her a $25\%$ discount. Pam buys 5 tickets using a coupon that gives her a $30\%$ discount. How many more dollars does Pam pay than Susan?


 $\mathrm{(A)}\ 2\qquad \mathrm{(B)}\ 5\qquad \mathrm{(C)}\ 10\qquad \mathrm{(D)}\ 15\qquad \mathrm{(E)}\ 20$


 
 \textbf{Answer: }C



\textbf{Solution}

$\textbf{Answer: \boxed{(C)}}$ Susan pays $(4)(0.75)(20) = 60$ dollars. Pam pays $(5)(0.70)(20) = 70$ dollars, so she pays $70-60=10$ more dollars than Susan.

 


	\item (\textit{AMC 10 10A 2007}) Define $a@b = ab - b^{2}$ and $a\#b = a + b - ab^{2}$. What is $\frac {6@2}{6\#2}$?


 $\text{(A)}\ - \frac {1}{2}\qquad \text{(B)}\ - \frac {1}{4}\qquad \text{(C)}\ \frac {1}{8}\qquad \text{(D)}\ \frac {1}{4}\qquad \text{(E)}\ \frac {1}{2}$


 
 \textbf{Answer: }A



\textbf{Solution}

$6@2$ must be equal to $6*2-2^2$ which is 8. $6\# 2$ is equal to $6+2-6*2^2$ which is $8-24 = -16$. Therefore $\frac{6@2}{6\# 2}$ must be equal to $\frac{8}{-16} = -\frac{1}{2}$. Therefore the solution is $\boxed{A}$.

 


	\item (\textit{AMC 10 10A 2007}) An aquarium has a rectangular base that measures 100 cm by 40 cm and has a height of 50 cm. It is filled with water to a height of 40 cm. A brick with a rectangular base that measures 40 cm by 20 cm and a height of 10 cm is placed in the aquarium. By how many centimeters does the water rise? 


 $\text{(A)}\ 0.5 \qquad \text{(B)}\ 1 \qquad \text{(C)}\ 1.5 \qquad \text{(D)}\ 2 \qquad \text{(E)}\ 2.5$


 
 \textbf{Answer: }D



\textbf{Solution}

The volume of the brick is $40 \times 20 \times 10 = 8000$. Thus the water volume rose $8000 = 100 \times 40 \times h \Longrightarrow h = 2\ \boxed{\mathrm{(D)}}$.

 


	\item (\textit{AMC 10 10A 2007}) The larger of two consecutive odd integers is three times the smaller. What is their sum?


 $\text{(A)}\ 4 \qquad \text{(B)}\ 8 \qquad \text{(C)}\ 12 \qquad \text{(D)}\ 16 \qquad \text{(E)}\ 20$


 
 \textbf{Answer: }A



\textbf{Solution}

Let the two consecutive odd integers be $a$, $a+2$. Then $a+2 = 3a$, so $a = 1, a + 2 = 3$ and their sum is $4\ \boxed{\mathrm{(A)}}$.

 


	\item (\textit{AMC 10 10A 2007}) A school store sells 7 pencils and 8 notebooks for $\$4.15$. It also sells 5 pencils and 3 notebooks for $\$1.77$. How much do 16 pencils and 10 notebooks cost?


 $\text{(A)}\ \$1.76 \qquad \text{(B)}\ \$5.84 \qquad \text{(C)}\ \$6.00 \qquad \text{(D)}\ \$6.16 \qquad \text{(E)}\ \$6.32$


 
 \textbf{Answer: }B



\textbf{Solution}

Let $p$ be cost of one pencil in dollars and $n$ the cost of one notebook in dollars. Then

 
 \begin{align*} 7p + 8n = 4.15 &\implies  35p + 40n = 20.75\\ 5p + 3n = 1.77 &\implies  35p + 21n = 12.39 \end{align*}


 Subtracting these equations yields $19n = 8.36 \Longrightarrow n = 0.44$. Solving backwards gives $p = 0.09$. Thus the answer is $16p + 10n = \text{(B)}\mathdollar 5.84$.

 



\textbf{Solution 2}

Since 5 pencils and 3 notebooks cost 1.77 dollars, then 3 times that or 15 pencils and 9 notebooks costs 5.31 dollars which is 1 pencil and 1 notebook off. Looking at answer choices, it can only be 5.84 so $\mathrm{(B)}$ .

 Note: 6.00 dollars would imply that 1 pencil and 1 notebook would cost more than 30\% of 5 pencils and 3 notebooks, which is incorrect. 

 


	\item (\textit{AMC 10 10A 2007}) At Euclid High School, the number of students taking the AMC 10 was $60$ in 2002, $66$ in 2003, $70$ in 2004, $76$ in 2005, $78$ in 2006, and is $85$ in 2007. Between what two consecutive years was there the largest percentage increase?


 $\text{(A)}\ 2002\ \text{and}\ 2003 \qquad \text{(B)}\ 2003\ \text{and}\ 2004 \qquad \text{(C)}\ 2004\ \text{and}\ 2005 \qquad \text{(D)}\ 2005\ \text{and}\ 2006 \qquad \text{(E)}\ 2006\ \text{and}\ 2007$


 
 \textbf{Answer: }A



\textbf{Solution 1}

We compute the percentage increases:

 
            \begin{enumerate}
                \item $\frac{66 - 60}{60} = 10\%$\item $\frac{70 - 66}{66} \approx 6\%$\item $\frac{76-70}{70} \approx 8.6\%$\item $\frac{78-76}{76} \approx 2.6\%$\item $\frac{85-78}{78} \approx 9\%$
            \end{enumerate}
            The answer is $\mathrm{(A)}$. 

 In fact, the answer follows directly from examining the differences between each year. The largest differences are $6$ and $7$. Due to the decreased starting number of students between $2002$ and $2003$, that interval will be our answer.

 Note: After we find the initial $10$\% difference to be our largest percent, we can increase each previous number by $10$\% and estimate whether that number is larger than the difference between the two numbers. ~Thyauraisoverwhelming

 



\textbf{Solution 2}

We make the numerator 1 and compare the denominators.

 
            \begin{enumerate}
                \item $\frac{66 - 60}{60} = \frac{1}{10}$
            \end{enumerate}
            
            \begin{enumerate}
                \item $\frac{70 - 66}{66} = \frac{1}{16.5}$
            \end{enumerate}
            
            \begin{enumerate}
                \item $\frac{76-70}{70} \approx \frac{1}{11.7}$
            \end{enumerate}
            
            \begin{enumerate}
                \item $\frac{78-76}{76} = \frac{1}{38}$
            \end{enumerate}
            
            \begin{enumerate}
                \item $\frac{85-78}{78} \approx \frac{1}{11.1}$
            \end{enumerate}
            
The answer is $\mathrm{(A)}$.

 


	\item (\textit{AMC 10 10A 2007}) Last year Mr. Jon Q. Public received an inheritance. He paid $20\%$ in federal taxes on the inheritance, and paid $10\%$ of what he had left in state taxes. He paid a total of $\$10500$ for both taxes. How many dollars was his inheritance?


 $(\mathrm {A})\ 30000 \qquad (\mathrm {B})\ 32500 \qquad(\mathrm {C})\ 35000 \qquad(\mathrm {D})\ 37500 \qquad(\mathrm {E})\ 40000$


 
 \textbf{Answer: }D



\textbf{Solution 1}

After paying his taxes, he has $0.8*0.9=0.72$ of his earnings left. Since $10500$ is $0.28$ of his income, he got a total of $\frac{10500}{0.28}=37500\  \mathrm{(D)}$.

 



\textbf{Solution 2}

Let his total inheritance be a number $x$. The amount of tax he has to pay at first can be represented as $0.2x$. The amount of money he has left over after this expense can be written as $0.8x$. He then has to pay $10\%$ tax on this remaining amount, which can be written as $0.1(0.8x) = 0.08x$. The combined expense of these taxes is $10500$, so $0.2x + 0.08x = 10500 \Longrightarrow 0.28x = 10500$. Therefore $x = \frac{10500}{0.28}\ = 37500  \mathrm{(D)}$. \href{https://artofproblemsolving.com/wiki/index.php?title=User:Ankitamc&action=edit&redlink=1}{Ankitamc} (\href{https://artofproblemsolving.com/wiki/index.php?title=User\_talk:Ankitamc&action=edit&redlink=1}{talk}) 13:32, 17 January 2021 (EST)AnkitAmc.~edited by mobius247

 


	\item (\textit{AMC 10 10A 2007}) Triangles $ABC$ and $ADC$ are isosceles with $AB=BC$ and $AD=DC$. Point $D$ is inside triangle $ABC$, angle $ABC$ measures 40 degrees, and angle $ADC$ measures 140 degrees. What is the degree measure of angle $BAD$?


 $\mathrm{(A)}\ 20\qquad \mathrm{(B)}\ 30\qquad \mathrm{(C)}\ 40\qquad \mathrm{(D)}\ 50\qquad \mathrm{(E)}\ 60$


 
 \textbf{Answer: }D



\textbf{Solution 1}


                    \begin{figure}[H]
                        \centering
                        \includegraphics[width=0.5\textwidth]{Solution/AMC_10_10A_2007_8_Solution_Figure_1.png}
                    \end{figure}
                    

 We angle chase and find out that:

 
                \begin{itemize}
                    \item $\angle DAC=\frac{180-140}{2} = 20$\item $\angle BAC=\frac{180-40}{2} = 70$\item $\angle BAD=\angle BAC- \angle DAC=50\ \mathrm{(D)}$
                \end{itemize}


                



\textbf{Solution 2}


                    \begin{figure}[H]
                        \centering
                        \includegraphics[width=0.5\textwidth]{Solution/AMC_10_10A_2007_8_Solution_Figure_2.png}
                    \end{figure}
                    

 Since triangle $ABC$ is isosceles we know that angle $\angle BAC = \angle BCA$.

 Also since triangle $ADC$ is isosceles we know that $\angle DAC = \angle DCA$.

 This implies that  $\angle BAD = \angle BCD$.

 Then the sum of the interior angles of quadrilateral $ABCD$ is  $40 + 220 + 2\angle BAD = 360$.

 Solving the equation, we get $\angle BAD = 50$.

 Therefore the answer is $\mathrm{(D)}$.

 



\textbf{Solution 3}

Using the previous drawing we know that ACD and DAC are both equal to 20 (40/2 or (180-40)/2). We also know that BAC and BCA are both 70 or (180-40)/2. Thus BDA is 70-20 which is equal to 50. 

 -RealityWrites

 


	\item (\textit{AMC 10 10A 2007}) Real numbers $a$ and $b$ satisfy the equations $3^{a} = 81^{b + 2}$ and $125^{b} = 5^{a - 3}$. What is $ab$?


 $\text{(A)}\ -60 \qquad \text{(B)}\ -17 \qquad \text{(C)}\ 9 \qquad \text{(D)}\ 12 \qquad \text{(E)}\ 60$


 
 \textbf{Answer: }E



\textbf{Solution 1}


 \[81^{b+2} = 3^{4(b+2)} = 3^a \Longrightarrow a = 4b+8\]

 And

 
 \[125^{b} = 5^{3b} = 5^{a-3} \Longrightarrow a - 3 = 3b\]

 Substitution gives $4b+8 - 3 = 3b \Longrightarrow b = -5$, and solving for $a$ yields $-12$. Thus $ab = 60\ \mathrm{(E)}$.

 



\textbf{Solution 1 another similar way}

Simplify equation $1$, which is $3^a=81^{b+2}$, to $3^a=3^{4b+8}$. 

 And

 Simplify equation $2$, which is $125^b=5^{a-3}$, to $5^{3b}=5^{a-3}$.

 Now, eliminate the bases from the simplified equations $1$ and $2$ to arrive at $a=4b+8$ and $3b=a-3$. Rewrite equation $2$ so that it is in terms of $a$. That would be $a=3b+3$.

 Since both equations are equal to $a$, and the values for $a$ and $b$ are constant for both equations, set the equations equal to each other.$4b+8=3b+3 \Longrightarrow b=-5$

 Now plug $b$, which is $-5$, back into one of the two earlier equations.$4(-5)+8=a \Longrightarrow -20+8=a \Longrightarrow a=-12$

 $(-12)(-5)=60$

 Therefore the correct answer is $\mathrm{(E)}$

 


	\item (\textit{AMC 10 10A 2007}) The Dunbar family consists of a mother, a father, and some children. The average age of the members of the family is $20$, the father is $48$ years old, and the average age of the mother and children is $16$. How many children are in the family?


 $\text{(A)}\ 2 \qquad \text{(B)}\ 3 \qquad \text{(C)}\ 4 \qquad \text{(D)}\ 5 \qquad \text{(E)}\ 6$


 
 \textbf{Answer: }E



\textbf{Solution 1}

Let $n$ be the number of children. Then the total ages of the family is $48 + 16(n+1)$, and the total number of people in the family is $n+2$. So

 
 \[20 = \frac{48 + 16(n+1)}{n+2} \Longrightarrow 20n + 40 = 16n + 64 \Longrightarrow n = 6\ \mathrm{(E)}.\]

 



\textbf{Solution 2}

Let $x$ be the number of the children and the mom. The father, who is $48$, plus the sum of the ages of the kids and mom divided by the number of kids and mom plus $1$ (for the dad) = $20$. This is because the average age of the entire family is $20.$This statement, written as an equation, is: 
 \[\frac{48+16x}{x+1}=20\]
 \[48+16x=20x+20\]
 \[4x=28\]
 \[x=7\]

 $7$ people - $1$ mom = $6$ children.

 Therefore, the answer is $\boxed{E}$

 



\textbf{Solution 3}

Let $m$ be the Mom's age.

 Let the number of children be $x$ and their average age be $y$. Their age totaled up is simply $xy$. 

 We have the following two equations:

 $\frac{m+48+xy}{2+x}=20$, where $m+48+xy$ is the family's total age and $2+x$ is the total number of people in the family.

 $\frac{m+48+xy}{2+x}=20$ 

 $m+48+xy=40+20x$

 The next equation is $\frac{m+xy}{1+x}=16$, where $m+xy$ is the total age of the Mom and the children, and $1+x$ is the number of children along with the Mom.

 $\frac{m+xy}{1+x}=16$

 $m+xy=16+16x$.

 We know the value for $m+xy$, so we substitute the value back in the first equation.

 $m+48+xy=40+20x$ 

 $(16+16x)+48=40+20x$.

 $x=6$. 

 Earlier, we set $x$ to be the number of children. Therefore, there are $\boxed{\text{(E)}  6}$ children.

 


	\item (\textit{AMC 10 10A 2007}) The numbers from $1$ to $8$ are placed at the vertices of a cube in such a manner that the sum of the four numbers on each face is the same. What is this common sum?


 $\text{(A)}\ 14 \qquad \text{(B)}\ 16 \qquad \text{(C)}\ 18 \qquad \text{(D)}\ 20 \qquad \text{(E)}\ 24$


 
 \textbf{Answer: }C



\textbf{Solution 1}

The sum of the numbers on one face of the cube is equal to the sum of the numbers on the opposite face of the cube; these $8$ numbers represent all of the vertices of the cube. Thus the answer is $\frac{1 + 2 + \cdots + 8}{2} = 18\ \mathrm{(C)}$.

 



\textbf{Solution 2}

Consider a number on a vertex. It will be counted in 3 different faces, since any vertex is the intersection of three edges. Therefore, each number $1,2,\cdots,7,8$ will be added into the total sum $3$ times. Therefore, our total sum is $3(1+2+\cdots+8)=108.$ Finally, since there are $6$ faces, our common sum is $\dfrac{108}{6} = 18\mathrm{(C) }.$

 


	\item (\textit{AMC 10 10A 2007}) Two tour guides are leading six tourists. The guides decide to split up. Each tourist must choose one of the guides, but with the stipulation that each guide must take at least one tourist. How many different groupings of guides and tourists are possible?


 $\text{(A)}\ 56 \qquad \text{(B)}\ 58 \qquad \text{(C)}\ 60 \qquad \text{(D)}\ 62 \qquad \text{(E)}\ 64$


 
 \textbf{Answer: }D



\textbf{Solution}

Each tourist has to pick in between the $2$ guides, so for $6$ tourists there are $2^6$ possible groupings. However, since each guide must take at least one tourist, we subtract the $2$ cases where a guide has no tourist. Thus the answer is $2^6 - 2 = \boxed{62}\ \mathrm{(D)}$.

 



\textbf{Solution 2}

Without loss of generality, let's call one of the tour guides tour guide A, and the other tour guide B. To count the number of total groupings of guides and tourists possible, we can count the number of ways some number of tourists go to tour guide A. Thus, we can see that the total number of groupings is:
 \[\binom{6}{1} + \binom{6}{2} + \binom{6}{3} + \binom{6}{4} + \binom{6}{5} = \boxed{62} \text{ (D)}\]

 



\textbf{Solution 3 (More detailed)}

Very similar to the solution above, we begin by naming the tour guides \href{https://artofproblemsolving.com/wiki/index.php/Without\_loss\_of\_generality}{/ WLOG}, tour guide $A$ and tour guide $B$. The possible groupings of tourists and tour guides are: Tour guide $A$ takes 1, Tour guide $B$ takes 5, where there are $\binom{6}{1}$ = $6$ ways to arrange the tourists.Tour guide $A$ takes 2, Tour guide $B$ takes 4, where there are $\binom{6}{2}$ = $15$ ways to arrange the tourists.Tour guide $A$ takes 3, Tour guide $B$ takes 3, where there are $\binom{6}{3}$ = $20$ ways to arrange the tourists.Tour guide $A$ takes 4, Tour guide $B$ takes 2, where there are $\binom{6}{2}$ = $15$ ways to arrange the tourists.Tour guide $A$ takes 5, Tour guide $B$ takes 1, where there are $\binom{6}{1}$ = $6$ ways to arrange the tourists.

 
Thus, the total number of different groupings of guides and tourists is equal to $6 + 15 + 20 + 15 + 6$, which equals $\boxed{\text{ (D)} 62}$

 


	\item (\textit{AMC 10 10A 2007}) Yan is somewhere between his home and the stadium. To get to the stadium he can walk directly to the stadium, or else he can walk home and then ride his bicycle to the stadium. He rides 7 times as fast as he walks, and both choices require the same amount of time. What is the ratio of Yan's distance from his home to his distance from the stadium?


 $\mathrm{(A)}\ \frac 23\qquad \mathrm{(B)}\ \frac 34\qquad \mathrm{(C)}\ \frac 45\qquad \mathrm{(D)}\ \frac 56\qquad \mathrm{(E)}\ \frac 78$


 
 \textbf{Answer: }B



\textbf{Solution 1}

Let the distance from Yan's initial position to the stadium be $a$ and the distance from Yan's initial position to home be $b$. We are trying to find $b/a$, and we have the following identity given by the problem:

 
 \begin{align*}a &= b + \frac{a+b}{7}\\ \frac{6a}{7} &= \frac{8b}{7} \Longrightarrow 6a = 8b\end{align*}


 Thus $b/a = 6/8 = 3/4$ and the answer is $\mathrm{(B)}\ \frac{3}{4}$

 Note: The identity can be given from an application of $t=\frac dr$.

 



\textbf{Solution 2}

Another way of solving this problem is by setting the distance between Yan's home and the stadium, thus filling in one variable.Let us set the distance between the two places to be $42z$, where $z$ is a random measurement (cause life, why not?). The time going to his home then riding his bike, which is $7$ times faster, to the stadium is equal to the time it takes for him to walk to the stadium. So the equation would be:Let $x=$ the distance from Yan's position to his home. Let $42z=$ the distance from Yan's home to the stadium.

 $42z-x=x+\frac{42z}{7}$

 $42z-x=x+6z$

 $36z=2x$

 $x=18z$

 But we're still not done with the question. We know that Yan is $18z$ from his home, and is also $42z-18z$ or $24z$ from the stadium.$\frac{18z}{24z}$  , the $z$'s cancel out, and we are left with $\frac{3}{4}$.Thus, the answer is $\mathrm{(B)}\ \frac{3}{4}$

 



\textbf{Solution 3}

Assume that the distance from home to the stadium is 1, and the distance from Yan to home is $b$. Also assume that the speed of walking is 1, so the speed of biking is 7. Thus$1-b=b+\frac{1}{7} \Longrightarrow b=\frac{3}{7}$.We want to find $\frac{b}{1-b}$ which is equal to

 $\frac{3}{7} \div \frac{4}{7} = \mathrm{(B)}\ \frac{3}{4}$

 



\textbf{Solution 4}


                    \begin{figure}[H]
                        \centering
                        \includegraphics[width=0.5\textwidth]{Solution/AMC_10_10A_2007_13_Solution_Figure_1.png}
                    \end{figure}
                    Let $H$ represent Yan's home, $S$ represent the stadium, and $Y$ represent Yan's current position. If Yan walks directly to the stadium, he will reach Point $P$ the same time he will reach $H$ if he is walking home. Since he bikes $7$ times as fast as he walks and the time is the same, the distance from his home to the stadium must be $7$ times the distance from $P$ to the stadium. If $PS=x$, then $HS=7x$ and $HP=6x$. Since Y is the midpoint of $\overline{HP}$, $HY=YP=3x$. Therefore, the ratio of Yan's distance from his home to his distance from the stadium is $\frac{YH}{YS}=\frac{3x}{4x}=\boxed{\mathrm{(B)}\ \frac{3}{4}}$

 


	\item (\textit{AMC 10 10A 2007}) A triangle with side lengths in the ratio $3 : 4 : 5$ is inscribed in a circle with radius $3$. What is the area of the triangle?


 $\mathrm{(A)}\ 8.64\qquad \mathrm{(B)}\ 12\qquad \mathrm{(C)}\ 5\pi\qquad \mathrm{(D)}\ 17.28\qquad \mathrm{(E)}\ 18$


 
 \textbf{Answer: }A



\textbf{Solution}


                    \begin{figure}[H]
                        \centering
                        \includegraphics[width=0.5\textwidth]{Solution/AMC_10_10A_2007_14_Solution_Figure_1.png}
                    \end{figure}
                    

 Since 3-4-5 is a \href{https://artofproblemsolving.com/wiki/index.php?title=Pythagorean\_triple}{Pythagorean triple}, the triangle is a \href{https://artofproblemsolving.com/wiki/index.php?title=Right\_triangle}{right triangle}. Since the hypotenuse is a \href{https://artofproblemsolving.com/wiki/index.php?title=Diameter}{diameter} of the \href{https://artofproblemsolving.com/wiki/index.php?title=Circumcircle}{circumcircle}, the hypotenuse is $2r = 6$. Then the other legs are $\frac{24}5=4.8$ and $\frac{18}5=3.6$. The area is $\frac{4.8 \cdot 3.6}2 = 8.64\ \mathrm{(A)}$

 



\textbf{Solution 2}

The hypotenuse of the triangle is a diameter of the circumcircle, so it has length $2 \cdot 3 = 6$. The triangle is similar to a 3-4-5 triangle with the ratio of their side lengths equal to $\frac{6}{5}$. The area of a 3-4-5 triangle is $\frac{3\cdot 4}{2} = 6$.The square of the ratio of their side lengths is equal to the ratio of their areas. Call the area of the triangle $A$. Therefore, $\left(\frac{6}{5}\right)^2 = \frac{A}{6} \Longrightarrow \frac{36}{25} = \frac{A}{6} \Longrightarrow A = \frac{36\cdot6}{25} = 8.64\ \mathrm{(A)}$

 



\textbf{Solution 3 (dimensional analysis)}

Using the theorem that the hypotenuse of a right triangle inscribed in a circle is the diameter, one can determine that the diameter of the triangle is 6. As a result the ratio between the triangular ratios and their real values is (6/5). Using dimensional analysis, one can see that because the ratio between lengths is (6/5), the ratio between areas is (6/5)\textasciicircum{}2. This gives ((3 * 4)(6/5)\textasciicircum{}2)/2 as the area of the triangle or answer choice (A).

 



\textbf{Solution 4}

Let the sides of the triangle be $3x, 4x, 5x$ for some $x$. Now use the fact that the area of a triangle is equal to $\frac{abc}{4R}$, where $a, b, c$ are side lengths and $R$ is the circumradius. Also note that due to the 3, 4, 5 triple, the triangle is right. Thus we get $\frac{12x^2}{2} = \frac{60x^3}{4\cdot 3}\Rightarrow x = \frac{6}{5} = 1.2$. Then the area is $12(1.2)^2/2 = \boxed{8.64 \text{ (A)}}$.

 


	\item (\textit{AMC 10 10A 2007}) Four circles of radius $1$ are each tangent to two sides of a square and externally tangent to a circle of radius $2$, as shown. What is the area of the square?


 

                    \begin{figure}[H]
                        \centering
                        \includegraphics[width=0.5\textwidth]{Problem/AMC_10_10A_2007_15_Problem_Figure_1.png}
                    \end{figure}
                    


 


 $\text{(A)}\ 32 \qquad \text{(B)}\ 22 + 12\sqrt {2}\qquad \text{(C)}\ 16 + 16\sqrt {3}\qquad \text{(D)}\ 48 \qquad \text{(E)}\ 36 + 16\sqrt {2}$


 
 \textbf{Answer: }B



\textbf{Solution 1}

Draw a square connecting the centers of the four small circles of radius $1$. This square has a diagonal of length $6$, as it includes the diameter of the big circle of radius $2$ and two radii of the small circles of radius $1$. Therefore, the side length of this square is 
 \[\frac{6}{\sqrt{2}} = 3\sqrt{2}.\] The side length of the larger square is $2$ units greater than the one found by connecting the midpoints, so its side length is 
 \[2 + 3\sqrt{2}.\] The area of the larger square is $(2+3\sqrt{2})^2 = 22 + 12\sqrt{2}$ $(B).$

 



\textbf{Solution 2}

We draw the diagonal of the square. This diagonal yields $2\sqrt{2}+1+1+2+2=2\sqrt{2}+6$. We know that the side length $s$ in terms of the diagonal $d$ is $s=\frac{d}{\sqrt{2}}$, so our side length is $\frac{2\sqrt{2}+6}{\sqrt{2}}$. However, we are trying to look for the area, so squaring $\frac{2\sqrt{2}+6}{\sqrt{2}}$ yields $\frac{44+24\sqrt{2}}{2}=\boxed{\text{(B)}22+12\sqrt{2}}$

 



\textbf{Solution 3}

Connect the points from centre of big circle to the top 2 circles. This creates a right isosceles triangle of length $2+1=3$. Thus its hypotenuse is $3\sqrt{2}$. Adding twice the radius of the top circle gets $2+3\sqrt{2}$. This is the side length of the square, thus its area is $\boxed{\text{(B)}22+12\sqrt{2}}$

 I'm bad at explaining things I'm sure its understandable with diagram

 


	\item (\textit{AMC 10 10A 2007}) Integers $a, b, c,$ and $d$, not necessarily distinct, are chosen independently and at random from 0 to 2007, inclusive. What is the probability that $ad-bc$ is even?


 $\mathrm{(A)}\ \frac 38\qquad \mathrm{(B)}\ \frac 7{16}\qquad \mathrm{(C)}\ \frac 12\qquad \mathrm{(D)}\ \frac 9{16}\qquad \mathrm{(E)}\ \frac 58$


 
 \textbf{Answer: }E



\textbf{Solution 1}

The only time when $ad-bc$ is even is when $ad$ and $bc$ are of the same \href{https://artofproblemsolving.com/wiki/index.php?title=Parity}{parity}. The chance of $ad$ being odd is $\frac 12 \cdot \frac 12 = \frac 14$, since the only way to have $ad$ be odd is to have both $a$ and $d$ be odd. As a result, $ad$ has a $\frac 34$ probability of being even. $bc$ also has a $\frac 14$ chance of being odd and a $\frac34$ chance of being even. Therefore, the probability that $ad-bc$ will be even is $\left(\frac 14\right)^2+\left(\frac 34\right)^2=\boxed {\mathrm{(E )} \frac 58\ }$.

 



\textbf{Solution 2 (casework)}

If we don't know our parity rules, we can check and see that $ad-bc$ is only even when $ad$ and $bc$ are of the same \href{https://artofproblemsolving.com/wiki/index.php?title=Parity}{parity} (as stated above). From here, we have two cases.

 Case 1: $odd-odd$ (which must be $o \cdot o-o \cdot o$). The probability for this to occur is $\left(\frac 12\right)^4 = \frac 1{16}$, because each integer has a $\frac 12$ chance of being odd.

 Case 2: $even-even$ (which occurs in 4 cases: $(e \cdot e-e \cdot e$), ($o \cdot e-o \cdot e$) (alternating of any kind), and ($e \cdot e-o \cdot e$) with its reverse, ($o \cdot e-e \cdot e$).

 Our first subcase of case 2 has a chance of $\frac 1{16}$ (same reasoning as above).

 Our second subcase of case 2 has a $\frac 14$ chance, since only the 2nd and 4th flip matter (or 1st and 3rd).

 Our third subcase of case 2 has a $\frac 18$ chance, because the 1st, 2nd, and either 3rd or 4th flip matter.

 Our fourth subcase of case 2 has a $\frac 18$ chance, because it's the same, just reversed.

 We sum these, and get our answer of $\frac 1{16} + \frac 1{16} + \frac 14 + \frac 18 + \frac 18 = \boxed {\mathrm{(E )} \frac 58\ }.$

 


	\item (\textit{AMC 10 10A 2007}) Suppose that $m$ and $n$ are positive integers such that $75m = n^{3}$. What is the minimum possible value of $m + n$?


 $\text{(A)}\ 15 \qquad \text{(B)}\ 30 \qquad \text{(C)}\ 50 \qquad \text{(D)}\ 60 \qquad \text{(E)}\ 5700$


 
 \textbf{Answer: }D



\textbf{Solution}

$3 \cdot 5^2m$ must be a perfect cube, so each power of a prime in the factorization for $3 \cdot 5^2m$ must be divisible by $3$. Thus the minimum value of $m$ is $3^2 \cdot 5 = 45$, which makes $n = \sqrt[3]{3^3 \cdot 5^3} = 15$. The minimum possible value for the sum of $m$ and $n$ is $\boxed {(D)60}.$

 



\textbf{Solution 2}

First, we need to prime factorize $75$. $75$ = $5^2 \cdot 3$. We need $75m$ to be in the form $x^3y^3$. Therefore, the smallest $m$ is $5 \cdot 3^2$. $m$ = 45, and since $5^3 \cdot 3^3 = 15^3$, our answer is $45 + 15$ = $\boxed {(D)60}$

 


	\item (\textit{AMC 10 10A 2007}) Consider the $12$-sided polygon $ABCDEFGHIJKL$, as shown. Each of its sides has length $4$, and each two consecutive sides form a right angle. Suppose that $\overline{AG}$ and $\overline{CH}$ meet at $M$. What is the area of quadrilateral $ABCM$?


 

                    \begin{figure}[H]
                        \centering
                        \includegraphics[width=0.5\textwidth]{Problem/AMC_10_10A_2007_18_Problem_Figure_1.png}
                    \end{figure}
                    


 


 $\text{(A)}\ \frac {44}{3}\qquad \text{(B)}\ 16 \qquad \text{(C)}\ \frac {88}{5}\qquad \text{(D)}\ 20 \qquad \text{(E)}\ \frac {62}{3}$


 
 \textbf{Answer: }C



\textbf{Solution 1}

We can obtain the solution by calculating the area of rectangle $ABGH$ minus the combined area of triangles $\triangle AHG$ and $\triangle CGM$.

 We know that triangles $\triangle AMH$ and $\triangle GMC$ are similar because $\overline{AH} \parallel \overline{CG}$. Also, since $\frac{AH}{CG} = \frac{3}{2}$, the ratio of the distance from $M$ to $\overline{AH}$ to the distance from $M$ to $\overline{CG}$ is also $\frac{3}{2}$. Solving with the fact that the distance from $\overline{AH}$ to $\overline{CG}$ is 4, we see that the distance from $M$ to $\overline{CG}$ is $\frac{8}{5}$.

 The area of $\triangle AHG$ is simply $\frac{1}{2} \cdot 4 \cdot 12 = 24$, the area of $\triangle CGM$ is $\frac{1}{2} \cdot \frac{8}{5} \cdot 8 = \frac{32}{5}$, and the area of rectangle $ABGH$ is $4 \cdot 12 = 48$.

 Taking the area of rectangle $ABGH$ and subtracting the combined area of $\triangle AHG$ and $\triangle CGM$ yields $48 - \left(24 + \frac{32}{5}\right) = \boxed{\frac{88}{5}}\ \text{(C)}$.

 



\textbf{Solution 2}


                    \begin{figure}[H]
                        \centering
                        \includegraphics[width=0.5\textwidth]{Solution/AMC_10_10A_2007_18_Solution_Figure_1.png}
                    \end{figure}
                    

 Extend $AB$ and $CH$ and call their intersection $N$.

 The triangles $CBN$ and $CGH$ are clearly similar with ratio $1:2$, hence $BN=2$ and thus $AN=6$. The area of the triangle $BCN$ is $\frac{2\cdot 4}2 = 4$.

 The triangles $MAN$ and $MGH$ are similar as well, and we now know that the ratio of their dimensions is $AN:GH = 6:4 = 3:2$.

 Draw altitudes from $M$ onto $AN$ and $GH$, let their feet be $M_1$ and $M_2$. We get that $MM_1 : MM_2 = 3:2$. Hence $MM_1 = \frac 35 \cdot 12 = \frac {36}5$. (An alternate way is by seeing that the set-up AHGCM is similar to the 2 pole problem (\href{https://artofproblemsolving.comhttp://www.artofproblemsolving.com/wiki/index.php/1951\_AHSME\_Problems/Problem\_30}{http://www.artofproblemsolving.com/wiki/index.php/1951\_AHSME\_Problems/Problem\_30}). Therefore, $MM_2$ must be $\frac{1}{\frac{1}{8}+\frac{1}{12}} = \frac{24}{5}$, by the harmonic mean. Thus, $MM_1$ must be $\frac{36}{5}$.)   

 Then the area of $AMN$ is $\frac 12 \cdot AN \cdot MM_1 = \frac{108}5$, and the area of $ABCM$ can be obtained by subtracting the area of $BCN$, which is $4$. Hence the answer is $\frac{108}5 - 4 =  \boxed{\frac{88}5}$.

 



\textbf{Solution 3}


                    \begin{figure}[H]
                        \centering
                        \includegraphics[width=0.5\textwidth]{Solution/AMC_10_10A_2007_18_Solution_Figure_2.png}
                    \end{figure}
                    

 We can use coordinates to solve this. Let $H=(0,0).$ Thus, we have $A=(0,12), C=(4,8), G=(4,0).$ Therefore, $AG$ has equation $-3x+12=y$ and $HC$ has equation $2x=y.$ Solving, we have $M=(\frac {12}{5},\frac {24}{5}).$ Using the \href{https://artofproblemsolving.com/wiki/index.php?title=Shoelace\_Theorem}{Shoelace Theorem} (or you could connect $LC$ and solve for the resulting triangle + trapezoid areas), we find $[ABCM]=\boxed{\mathrm{(C) \ }\dfrac{88}{5}}.$

 



\textbf{Solution 4}

Let $J = (0,0)$ so we can find the equations of lines $AG$ and $HC$, and find the coordinates of point $M$. 

 We find the equation of line $AG$ to be $y = -3x + 20$ and line $HC$ to be $y = 2x + 12$. Setting these equations equal to each other, we solve for $M$:
 \[-3x + 20 = 2x + 12 \implies x = \frac{8}{5}\]Substituting $x$ back in gives $M = \left(\frac{32}{5}, \frac{4}{5}\right)$.

 We can also find the area of the smaller triangle in the top square by substituting $y = 4$ into our equations to see that the base segment $LC$ is split into two segments of length $\frac{4}{3}$ and $\frac{8}{3}$, respectively. 

 The area of this smaller triangle portion is:
 \[\frac{4 \times \frac{4}{3}}{2} = \frac{8}{3}\]

 Subtracting this from the total area of $16$, we find that the top portion of the quadrilateral's area is:
 \[16 - \frac{8}{3} = \frac{40}{3}\]

 For the bottom triangle portion with base $\frac{8}{3}$, the height is simply $4 - \frac{4}{5} = \frac{16}{5}$. The area of this triangle is:
 \[\frac{\frac{8}{3} \times \frac{16}{5}}{2} = \frac{64}{15}\]

 Finally, adding the two areas together gives:
 \[\frac{40}{3} + \frac{64}{15} = \frac{200}{15} + \frac{64}{15} = \frac{264}{15} = \frac{88}{5}\]

 Thus, our answer is $\boxed{\text{(C)}\ \frac{88}{5}}$.

 


	\item (\textit{AMC 10 10A 2007}) A paint brush is swept along both diagonals of a square to produce the symmetric painted area, as shown. Half the area of the square is painted. What is the ratio of the side length of the square to the brush width?


 

                    \begin{figure}[H]
                        \centering
                        \includegraphics[width=0.5\textwidth]{Problem/AMC_10_10A_2007_19_Problem_Figure_1.png}
                    \end{figure}
                    


 


 $\text{(A)}\ 2\sqrt {2} + 1 \qquad \text{(B)}\ 3\sqrt {2}\qquad \text{(C)}\ 2\sqrt {2} + 2 \qquad \text{(D)}\ 3\sqrt {2} + 1 \qquad \text{(E)}\ 3\sqrt {2} + 2$


 
 \textbf{Answer: }C



\textbf{Solution 1}

\href{https://artofproblemsolving.com/wiki/index.php?title=Without\_loss\_of\_generality}{Without loss of generality}, let the side length of the square be $1$ unit. The area of the painted area is $\frac{1}2$ of the area of the larger square, so the total unpainted area is also $\frac{1}{2}$.Each of the $4$ unpainted triangles has area $\frac{1}8$. It can be observed that these triangles are isosceles right triangles, so let $a$ be the side length of one of the smaller triangles:

 $\frac{a^2}2 = \frac{1}8 \rightarrow a= \frac{1}2$

 

 The hypotenuse of the triangle is $\frac{\sqrt{2}}2$. The corners of the painted area are also isosceles right triangles with side length $\frac{1-\frac{\sqrt{2}}2}2 = \frac{1}2-\frac{\sqrt2}4$. Its hypotenuse is equal to the width of the paint, and is $\frac{\sqrt{2}}2-\frac{1}2$. The answer we are looking for is thus $\frac{1}{\frac{\sqrt{2}}2-\frac{1}2}$. Multiply the numerator and the denominator by $\frac{\sqrt{2}}2+\frac{1}2$ to simplify, and you get $\frac{\frac{\sqrt{2}}{2}+\frac{1}{2}}{\frac{2}{4}-\frac{1}{4}}$ or $4\left(\frac{\sqrt{2}}{2}+\frac{1}{2}\right)$ which is $\boxed{(C)\ 2\sqrt{2}+2}$.

 


 



\textbf{Solution 2}

Again, have the length of the square equal to $1$ and let the width of each individual stripe be $n$. Note that you can split each stripe into one rectangle and two isosceles right triangles at the corners. Then the area of each stripe is $n(\sqrt{2}-\frac{n}{2})=\sqrt{2}n-\frac{n^2}{2}$. The area covered by the two total stripes is twice the area of one stripe, minus the area in the intersection of the stripes, which is a square with side length $n$. This area is equal to $\frac{1}{2}$. So:

 $2\sqrt2n-2n^2=\frac{1}{2}\rightarrow-2n^2+2\sqrt2n-\frac{1}{2}=0$.

 

 By the \href{https://artofproblemsolving.com/wiki/index.php?title=Quadratic\_formula}{quadratic formula}, 

 $n=\frac{-2\sqrt{2}\pm\sqrt{(2\sqrt{2})^2-4(-2)(-\frac{1}{2})}}{2(-2)}\rightarrow n=\frac{-2\sqrt{2}\pm 2}{-4}\rightarrow n=\frac{\sqrt{2}\pm 1}{2}$

 

 $\frac{\sqrt{2}+1}{2}$ is greater than the length of the square, so it is impossible for that to be the brush width. Thus $n=\frac{\sqrt{2}-1}{2}$. We want to find $\frac{1}{n}$, and $\frac{1}{n}=\frac{2}{\sqrt{2}-1}$. Multiply the numerator and the denominator by $\sqrt{2}+1$, 

 $\frac{1}{n}=\frac{2(\sqrt{2}+1)}{(\sqrt{2}-1)(\sqrt{2}+1)}=\boxed{2\sqrt{2}+2}\rightarrow C$

 

 


	\item (\textit{AMC 10 10A 2007}) Suppose that the number $a$ satisfies the equation $4 = a + a^{ - 1}$. What is the value of $a^{4} + a^{ - 4}$?


 $\text{(A)}\ 164 \qquad \text{(B)}\ 172 \qquad \text{(C)}\ 192 \qquad \text{(D)}\ 194 \qquad \text{(E)}\ 212$


 
 \textbf{Answer: }D



\textbf{Solution 1 (Decreases the Powers)}

Note that for all real numbers $k,$ we have $a^{2k} + a^{-2k} + 2 = (a^{k} + a^{-k})^2,$ from which 
 \[a^{2k} + a^{-2k} = (a^{k} + a^{-k})^2-2.\] We apply this result twice to get the answer:
 \begin{align*} a^4 + a^{-4} &= (a^2 + a^{-2})^2 - 2 \\ &= [(a + a^{-1})^2 - 2]^2 - 2 \\ &= \boxed{\textbf{(D)}\ 194}. \end{align*}
~Azjps ~MRENTHUSIASM

 



\textbf{Solution 2 (Increases the Powers)}

Squaring both sides of $a+a^{-1}=4$ gives $a^2+a^{-2}+2=16,$ from which $a^2+a^{-2}=14.$

 Squaring both sides of $a^2+a^{-2}=14$ gives $a^4+a^{-4}+2=196,$ from which $a^4+a^{-4}=\boxed{\textbf{(D)}\ 194}.$

 



\textbf{Solution 3 (Detailed Explanation of Solution 2)}

The detailed explanation of Solution 2 is as follows:
 \begin{alignat*}{8} a+a\textasciicircum{}{-1}&=4 \\ (a+a\textasciicircum{}{-1})\textasciicircum{}2&=4\textasciicircum{}2 \\ a\textasciicircum{}2+2aa\textasciicircum{}{-1}+a\textasciicircum{}{-2}&=16 \\ a\textasciicircum{}2+a\textasciicircum{}{-2}&=16-2&&=14 \\ (a\textasciicircum{}2+a\textasciicircum{}{-2})\textasciicircum{}2&=14\textasciicircum{}2 \\ a\textasciicircum{}4+2a\textasciicircum{}2a\textasciicircum{}{-2}+a\textasciicircum{}{-4}&=196 \\ a\textasciicircum{}4+a\textasciicircum{}{-4}&=196-2&&=\boxed{\textbf{(D)}\ 194}. \\ \end{alignat*}~MathFun1000 ~MRENTHUSIASM

 



\textbf{Solution 4 (Binomial Theorem)}

Squaring both sides of $a+a^{-1}=4$ gives $a^2+a^{-2}+2=16,$ from which $a^2+a^{-2}=14.$

 Applying the Binomial Theorem, we raise both sides of $a+a^{-1}=4$ to the fourth power:
 \begin{align*} \binom40a^4a^0+\binom41a^3a^{-1}+\binom42a^2a^{-2}+\binom43a^1a^{-3}+\binom44a^0a^{-4}&=256 \\ a^4+4a^2+6+4a^{-2}+a^{-4}&=256 \\ \left(a^4+a^{-4}\right)+4\left(a^2+a^{-2}\right)&=250 \\ \left(a^4+a^{-4}\right)+4(14)&=250 \\ a^4+a^{-4}&=\boxed{\textbf{(D)}\ 194}. \end{align*}
~MRENTHUSIASM

 



\textbf{Solution 5 (Solves for a)}

We multiply both sides of $4=a+a^{-1}$ by $a,$ then rearrange: 
 \[a^2-4a+1=0.\]

 We apply the Quadratic Formula to get $a=2\pm\sqrt3.$ 

 By Vieta's Formulas, note that the roots are reciprocals of each other. Therefore, both values of $a$ produce the same value of $a^4+a^{-4}:$
 \begin{align*} a^4+a^{-4}&=\left(2+\sqrt{3}\right)^4 + \left(2+\sqrt{3}\right)^{-4} \\ &=\left(2+\sqrt{3}\right)^4+\left(2-\sqrt{3}\right)^4 &&(*) \\ &=\boxed{\textbf{(D)}\ 194}. \end{align*}
\textbf{Remarks about $\boldsymbol{(*)}$}

 <ol style="margin-left: 1.5em;"> \item To find the fourth power of a sum/difference, we can first square that sum/difference, then square the result. 

 \item When we expand the fourth powers and combine like terms, the irrational terms will cancel.

 </ol>



\textbf{Solution 6 (Newton's Sums)}

From the first sentence of Solution 4, we conclude that $a$ and $a^{-1}$ are the roots of $x^2-4x+1=0.$ Let
 \begin{align*} P_1&=a+a^{-1}, \\ P_2&=a^2+a^{-2}, \\ P_3&=a^3+a^{-3}, \\ P_4&=a^4+a^{-4}. \end{align*}
By Newton's Sums, we have
 \begin{alignat*}{12} &1\cdot P\_1-4\cdot 1&&=0 &&\qquad\implies\qquad P\_1&&=4, \\ &1\cdot P\_2-4\cdot P\_1+1\cdot2 &&=0 &&\qquad\implies\qquad P\_2&&=14, \\ &1\cdot P\_3-4\cdot P\_2+1\cdot P\_1&&=0 &&\qquad\implies\qquad P\_3&&=52, \\ &1\cdot P\_4-4\cdot P\_3+1\cdot P\_2&&=0 &&\qquad\implies\qquad P\_4&&=\boxed{\textbf{(D)}\ 194}. \end{alignat*}~Albert1993 ~MRENTHUSIASM

 



\textbf{Solution 7 (Answer Choices)}

Note that 
 \[a^{4} + a^{-4} = (a^{2} + a^{-2})^{2} - 2.\] We guess that $a^{2} + a^{-2}$ is an integer, so the answer must be $2$ less than a perfect square. The only possibility is $\boxed{\textbf{(D)}\ 194}.$

 


	\item (\textit{AMC 10 10A 2007}) A sphere is inscribed in a cube that has a surface area of $24$ square meters. A second cube is then inscribed within the sphere. What is the surface area in square meters of the inner cube?


 $\text{(A)}\ 3 \qquad \text{(B)}\ 6 \qquad \text{(C)}\ 8 \qquad \text{(D)}\ 9 \qquad \text{(E)}\ 12$


 
 \textbf{Answer: }C



\textbf{Solution 1}


                    \begin{figure}[H]
                        \centering
                        \includegraphics[width=0.5\textwidth]{Solution/AMC_10_10A_2007_21_Solution_Figure_1.png}
                    \end{figure}
                    

 We rotate the smaller cube around the sphere such that two opposite vertices of the cube are on opposite faces of the larger cube. Thus the main diagonal of the smaller cube is the side length of the outer square.

 Let $S$ be the surface area of the inner square. The ratio of the areas of two similar figures is equal to the square of the ratio of their sides. As the diagonal of a cube has length $s\sqrt{3}$ where $s$ is a side of the cube, the ratio of a side of the inner square to a side of the outer square is $\frac{1}{\sqrt{3}}$  (since the side of the outer square = the diagonal of the inner square). So we have $\frac{S}{24} = \left(\frac{1}{\sqrt{3}}\right)^2$. Thus $S = 8\Rightarrow \mathrm{\boxed{(C)}}$.

 



\textbf{Solution 2 (computation)}

The area of each face of the outer cube is $\frac {24}{6} = 4$, and the edge length of the outer cube is $2$. This is also the \href{https://artofproblemsolving.com/wiki/index.php?title=Diameter}{diameter} of the sphere, and thus the length of a long diagonal of the inner cube.

 A long diagonal of a cube is the hypotenuse of a right triangle with a side of the cube and a face diagonal of the cube as legs.  If a side of the cube is $x$, we see that $2 = \sqrt {x^{2} + (\sqrt {2}x)^{2}}\Rightarrow x = \frac {2}{\sqrt {3}}$.

 Thus the surface area of the inner cube is $6x^{2} = 6\left(\frac {2}{\sqrt {3}}\right )^{2} = 8$.

 



\textbf{Solution 3, from AoPS}

Since the \href{https://artofproblemsolving.com/wiki/index.php?title=Surface\_area}{surface area} of the original \href{https://artofproblemsolving.com/wiki/index.php?title=Cube}{cube} is 24 square meters, each face of the cube has a surface area of $24/6 = 4$ square meters, and the side length of this cube is 2 meters. The sphere inscribed within the cube has diameter 2 meters, which is also the length of the diagonal of the cube inscribed in the sphere. Let $l$ represent the side length of the inscribed cube. Applying the \href{https://artofproblemsolving.com/wiki/index.php?title=Pythagorean\_Theorem}{Pythagorean Theorem} twice gives
 \[l^2 + l^2 + l^2 = 2^2 = 4.\]Hence each face has surface area
 \[l^2 = \frac{4}{3} \ \text{square meters}.\]So the surface area of the inscribed cube is $6\cdot \frac43 = \boxed{8}$ square meters.

 



\textbf{Solution 4}

First of all, it is pretty easy to see the length of each edge of the bigger cube is $2$ so the radius of the sphere is $1$. We know that when a cube is inscribed in a sphere, assuming the edge length of the square is $x$ the radius can be presented as $\frac {\sqrt3x}{2} = 1$, we can solve that $x=\frac {2\sqrt3}{3}$ and now we only need to apply the basic formula to find the surface and we got our final answer as $\frac {2\sqrt3}{3}*\frac {2\sqrt3}{3}*6=8$ and the final answer is $\boxed{C}$

 


	\item (\textit{AMC 10 10A 2007}) A finite sequence of three-digit integers has the property that the tens and units digits of each term are, respectively, the hundreds and tens digits of the next term, and the tens and units digits of the last term are, respectively, the hundreds and tens digits of the first term. For example, such a sequence might begin with the terms 247, 475, and 756 and end with the term 824. Let $S$ be the sum of all the terms in the sequence. What is the largest prime factor that always divides $S$? 


 $\mathrm{(A)}\ 3\qquad \mathrm{(B)}\ 7\qquad \mathrm{(C)}\ 13\qquad \mathrm{(D)}\ 37\qquad \mathrm{(E)}\ 43$


 
 \textbf{Answer: }D



\textbf{Solution}

A given digit appears as the hundreds digit, the tens digit, and the units digit of a term the same number of times. Let $k$ be the sum of the units digits in all the terms. Then $S=111k=3 \cdot 37k$, so $S$ must be divisible by $37\ \mathrm{(D)}$. To see that it need not be divisible by any larger prime, the sequence $123, 231, 312$ gives $S=666=2 \cdot 3^2 \cdot 37\Rightarrow \mathrm{\boxed{(D)~~37}}$.

 



\textbf{Solution 2 (Principle of Fractions)}

Let the first term be $100a+10b+c$. The next term, $100b+10c+d$, then $100c+10d+e$, then $100d+10e+f$, and so on until $100w+10x+y$ and then $100x+10y+z$. Then, the sum $S$ is just the first term $100a$, added with $110b$, and then added with $111(c+d+\dots+x)$, followed by $11y$ and finally $z$. In mathematical terms,
 \[S = 100a+110b+111(c+d+\dots+x) + 11y + z\].

 The largest prime $p~|~S$ satisfies the property that it must equivalate itself as an integer when divided and added to all other terms or must divide one of the coefficients. If the former is true, then $100+110+111+11+1 = 333$ must equivalate to zero modulo that prime, implying that 3 and 37 are possible equivalencies. To verify no other cases, we check the latter case. Out of all plausible prime divisors of each respective coefficient we see that the extremum $111$ contains the largest, which is 37. We check modulo 37 on $S$ but quickly arise the same conclusion as the prior case, supporting that 37 is indeed the largest prime.

 The requested largest prime is $\boxed{37}$. $\square$

 \textbf{Post Problem Note} Note that here, $y$ and $z$ are interchangeable for $a$ and $b$ given the problem directory, and by substitution one will arrive to the same conclusion as solution 1.

 



\textbf{Solution 3 (Using the provided sequence/Guessing)}

Since a prime number will always divide $S$, the prime number must work on any sequence that follows the stated rule. Since the solution kindly provides us with the terms $247,475,756,824$,all we have to do is finish the sequence.

 One way of finishing the sequence is by adding $568,682$ (before 824) to the sequence and adding them all together. After this, we just divide the result by each of the answer choices. Of the five, $37$ is the largest prime number that divides this sequence, so the answer is ${\boxed{(D)}}$.~Banspeedrun

 



\textbf{Solution 4 (a better way of doing solution 3)}

Because the question says which prime number will \textbf{ \textit{ ALWAYS }} divide $S$, the prime number would divide any sequence that follows the rule. So, to test this, create two sequences with three terms each- $111, 111, 111$ and $222, 222, 222$. The sum for both (respectively) is $333$ and $666$. Obviously, $3$ divides both. Now, test the other choices on both numbers, but don't if it doesn't divide one. An example is $\frac{333}{7}$. Because it doesn't divide, don't check $\frac{666}{7}$. Use the same logic for $\frac{333}{13}$. However, when you test $\frac{333}{37}$, it works. Now, you must check $\frac{666}{37}$ to ensure it isn't a special case. After doing this, you get that both are divisible. Now, check $\frac{333}{43}$ and see that it doesn't work. 

 
So the answer is ${\boxed{(D)}}$

 -N.

 


	\item (\textit{AMC 10 10A 2007}) How many ordered pairs $(m,n)$ of positive integers, with $m \ge n$, have the property that their squares differ by $96$? 


 $\text{(A)}\ 3 \qquad \text{(B)}\ 4 \qquad \text{(C)}\ 6 \qquad \text{(D)}\ 9 \qquad \text{(E)}\ 12$


 
 \textbf{Answer: }B



\textbf{Solution 1}


 \[m^2 - n^2 = (m+n)(m-n) = 96 = 2^{5} \cdot 3\] 

 For every two factors $xy = 96$, we have $m+n=x, m-n=y \Longrightarrow m = \frac{x+y}{2}, n = \frac{x-y}{2}$. It follows that the number of ordered pairs $(m,n)$ is given by the number of ordered pairs $(x,y): xy=96, x > y > 0$. There are $(5+1)(1+1) = 12$ factors of $96$, which give us six pairs $(x,y)$. However, since $m,n$ are positive integers, we also need that $\frac{x+y}{2}, \frac{x-y}{2}$ are positive integers, so $x$ and $y$ must have the same \href{https://artofproblemsolving.com/wiki/index.php?title=Parity}{parity}. Thus we exclude the factors $(x,y) = (1,96)(3,32)$, and we are left with $4$ pairs $\mathrm{(B)}$.

 



\textbf{Solution 2}

We first start as in Solution 1.However, as an alternative, we could also "give" each of the factors a factor of $2.$ This would force each one to be even. Now we have $\dfrac{96} {4} = 24,$ and since $24= 2^3 \cdot 3,$ the number of factors is $4 \cdot 2 = 8.$ We then divide by $2$ because $m-n \le m+n.$ This gives $4,$ as desired.

 



\textbf{Solution 3}

Similarly to the solution above, write $96$ as $2^5\cdot3^1$. To find the number of distinct factors, add $1$ to both exponents and multiply, which gives us $6\cdot2=12$ factors. Divide by $2$ since $m$ must be greater than or equal to $n$. We don't need to worry about $m$ and $n$ being equal because $96$ is not a perfect square. Finally, subtract the two cases above for the same reason to get $\mathrm{(B)}$.

 



\textbf{Solution 4}

Find all of the factor pairs of $96$: $(1,96),(2,48),(3,32),(4,24),(6,16),(8,12).$ You can eliminate $(1,96)$ and ($3,32)$ because you cannot have two numbers add to be an even number and have an odd difference at the same time without them being a decimal. You only have $4$ pairs left, so the answer is $\boxed{\textbf{(B)}\; 4}$.

 


 


	\item (\textit{AMC 10 10A 2007}) Circles centered at $A$ and $B$ each have radius $2$, as shown. Point $O$ is the midpoint of $\overline{AB}$, and $OA = 2\sqrt {2}$. Segments $OC$ and $OD$ are tangent to the circles centered at $A$ and $B$, respectively, and $EF$ is a common tangent. What is the area of the shaded region $ECODF$?


 

                    \begin{figure}[H]
                        \centering
                        \includegraphics[width=0.5\textwidth]{Problem/AMC_10_10A_2007_24_Problem_Figure_1.png}
                    \end{figure}
                    


 


 $\text{(A)}\ \frac {8\sqrt {2}}{3} \qquad \text{(B)}\ 8\sqrt {2} - 4 - \pi \qquad \text{(C)}\ 4\sqrt {2} \qquad \text{(D)}\ 4\sqrt {2} + \frac {\pi}{8} \qquad \text{(E)}\ 8\sqrt {2} - 2 - \frac {\pi}{2}$


 
 \textbf{Answer: }B



\textbf{Solution}

The area we are trying to find is simply $ABFE-(\overarc{AEC}+\triangle{ACO}+\triangle{BDO}+\overarc{BFD}).\overline{EF}\parallel\overline{AB}$.  Thus, $ABFE$ is a \href{https://artofproblemsolving.com/wiki/index.php?title=Rectangle}{rectangle}, and so its area is $b\times{h}=2\times{(AO+OB)}=2\times{2(2\sqrt{2})}=8\sqrt{2}$.

 Since $\overline{OC}$ is tangent to circle $A$, $\triangle{ACO}$ is a right triangle.  We know $AO=2\sqrt{2}$ and $AC=2$, so $\triangle{ACO}$ is an isosceles right triangle, and has $\overline{CO}$ with length $2$.  The area of $\triangle{ACO}=\frac{1}{2}bh=2$.  By symmetry, $\triangle{ACO}\cong\triangle{BDO}$, and so the area of $\triangle{BDO}$ is also $2$.

 $\overarc{AEC}$ (or $\overarc{BFD}$, for that matter) is $\frac{1}{8}$ the area of its circle since $\angle{OAC}$ is 45 degrees and $\angle{OAE}$ forms a right \href{https://artofproblemsolving.com/wiki/index.php?title=Triangle}{triangle}.  Thus $\overarc{AEC}$ and $\overarc{BFD}$ both have an area of $\frac{\pi}{2}$.

 Plugging all of these areas back into the original equation yields $8\sqrt{2}-(\frac{\pi}{2}+2+2+\frac{\pi}{2})=8\sqrt{2}-(4+\pi)=\boxed{(B)\ 8\sqrt{2}-4-\pi}$.

 


	\item (\textit{AMC 10 10A 2007}) For each positive integer $n$, let $S(n)$ denote the sum of the digits of $n.$ For how many values of $n$ is $n + S(n) + S(S(n)) = 2007?$


 $\mathrm{(A)}\ 1 \qquad \mathrm{(B)}\ 2 \qquad \mathrm{(C)}\ 3 \qquad \mathrm{(D)}\ 4 \qquad \mathrm{(E)}\ 5$


 
 \textbf{Answer: }D



\textbf{Solution 1}

For the sake of notation, let $T(n) = n + S(n) + S(S(n))$. Obviously $n<2007$. Then the maximum value of $S(n) + S(S(n))$ is when $n = 1999$, and the sum becomes $28 + 10 = 38$. So the minimum bound is $2007-38=1969$. We do \href{https://artofproblemsolving.com/wiki/index.php?title=Casework}{casework} upon the tens digit:

 Case 1: $196u \Longrightarrow u = 9$. Easy to directly disprove.

 Case 2: $197u$. $S(n) = 1 + 9 + 7 + u = 17 + u$, and $S(S(n)) = 8+u$ if $u \le 2$ and $S(S(n)) = 2 + (u-3) = u-1$ otherwise. 

 
            \begin{enumerate}[label=(\roman*)]
                \item Subcase a: $T = 1970 + u + 17 + u + 8 + u = 1995 + 3u = 2007 \Longrightarrow u = 4$. This exceeds our bounds, so no solution here.\item Subcase b: $T = 1970 + u + 17 + u + u - 1 = 1986 + 3u = 2007 \Longrightarrow u = 7$. First solution.
            \end{enumerate}
            Case 3: $198u$. $S(n) = 18 + u$, and $S(S(n)) = 9 + u$ if $u \le 1$ and $2 + (u-2) = u$ otherwise.

 
            \begin{enumerate}[label=(\roman*)]
                \item Subcase a: $T = 1980 + u + 18 + u + 9 + u = 2007 + 3u = 2007 \Longrightarrow u = 0$. Second solution.\item Subcase b: $T = 1980 + u + 18 + u + u = 1998 + 3u = 2007 \Longrightarrow u = 3$. Third solution.
            \end{enumerate}
            Case 4: $199u$. But $S(n) > 19$, and $n + S(n)$ clearly sum to $> 2007$.

 Case 5: $200u$. So $S(n) = 2 + u$ and $S(S(n)) = 2 + u$ (recall that $n < 2007$), and $2000 + u + 2 + u + 2 + u = 2004 + 3u = 2007 \Longrightarrow u = 1$. Fourth solution.

 In total we have $4 \mathrm{(D)}$ solutions, which are $1977, 1980, 1983,$ and $2001$.

 



\textbf{Solution 2}

Clearly, $n > 1950$. We can break this into three cases:

 Case 1: $n \geq 2000$

 
            \begin{enumerate}[label=(\roman*)]
                \item Inspection gives $n = 2001$.
            \end{enumerate}
            Case 2: $n < 2000$, $n = \overline{19xy}$, $x + y < 10$

 
            \begin{enumerate}[label=(\roman*)]
                \item If you set up an equation, it reduces to
            \end{enumerate}
            $4x + y = 32$

 
            \begin{enumerate}[label=(\roman*)]
                \item which has as its only solution satisfying the constraints $x = 8$, $y = 0$.
            \end{enumerate}
            Case 3: $n < 2000$, $n = \overline{19xy}$, $x + y \geq 10$

 
            \begin{enumerate}[label=(\roman*)]
                \item This reduces to
            \end{enumerate}
            
            \begin{enumerate}[label=(\roman*)]
                \item $4x + y = 35$. The only two solutions satisfying the constraints for this equation are $x = 7$, $y = 7$ and $x = 8$, $y = 3$.
            \end{enumerate}
            The solutions are thus $1977, 1980, 1983, 2001$ and the answer is $\mathrm{(D)}\  4$.

 



\textbf{Solution 3 (Fastest casework)}

It is well-known that $n \equiv S(n)\equiv S(S(n)) \pmod{9}.$ Substituting, we have that 
 \[n+n+n \equiv 2007 \pmod{9} \implies n \equiv 0 \pmod{3}.\] Since $n \leq 2007,$ we must have that $\max S(n)=1+9+9+9=28.$ Now, we list out the possible values for $S(n)$ in a table, noting that it is a multiple of $3$ because $n$ is a multiple of $3.$ 

 $\begin{tabular}{c|c c c c c c c c c c}   S(n) & 0 & 3 & 6 & 9 & 12 & 15 & 18 & 21 & 24 & 27 \\ \end{tabular}$

 

 Then, we compute the corresponding values of $S(S(n)).$

 $\begin{tabular}{c|c c c c c c c c c c}   S(n) & 0 & 3 & 6 & 9 & 12 & 15 & 18 & 21 & 24 & 27 \\ \hline S(S(n)) & 0 & 3 & 6 & 9 & 3 & 6 & 9 & 3 & 6 & 9 \\ \end{tabular}$

 

 Finally, we may compute the corresponding values of $n$ using the fact that $n=2007-S(n)-S(S(n)).$

 $\begin{tabular}{c|c c c c c c c c c c}  n & 2007 & 2001 & 1995 & 1989 & 1992 & 1986 & 1980 & 1983 & 1977 & 1971 \\ \hline S(n) & 0 & 3 & 6 & 9 & 12 & 15 & 18 & 21 & 24 & 27 \\ \hline S(S(n)) & 0 & 3 & 6 & 9 & 3 & 6 & 9 & 3 & 6 & 9 \\ \end{tabular}$

 

 Notice how all conditions are designed to be satisfied except whether $S(n)$ is accurate with respect to $n.$ So, the only thing that remains is to check this. We may eliminate, for example, when $n=2007$ we have $S(n)=9$ while the table states that it is $0.$ Proceeding similarly, we obtain the following table.

 $\begin{tabular}{c|c c c c c c c c c c}  n & \cancel{2007} & \boxed{2001} & \cancel{1995} & \cancel{1989} & \cancel{1992} & \cancel{1986} & \boxed{1980} & \boxed{1983} & \boxed{1977} & \cancel{1971} \\ \hline  S(n) & 0 & 3 & 6 & 9 & 12 & 15 & 18 & 21 & 24 & 27 \\ \hline S(S(n)) & 0 & 3 & 6 & 9 & 3 & 6 & 9 & 3 & 6 & 9 \\ \end{tabular}$

 

 It follows that there are $\boxed{4}\implies \textbf{(D)}$ possible values for $n.$ ~samrocksnature

 



\textbf{Solution 4}

As in Solution 1, we note that $S(n)\leq 28$ and $S(S(n))\leq 10$. 
Obviously, $n\equiv S(n)\equiv S(S(n)) \pmod 9$.  
As $2007\equiv 0 \pmod 9$, this means that $n\bmod 9 \in\{0,3,6\}$, or equivalently that $n\equiv S(n)\equiv S(S(n))\equiv 0 \pmod 3$.

 Thus $S(S(n))\in\{3,6,9\}$. For each possible $S(S(n))$ we get three possible $S(n)$.  
(E. g., if $S(S(n))=6$, then $S(n)=x$ is a number such that $x\leq 28$ and $S(x)=6$, therefore $S(n)\in\{6,15,24\}$.)

 For each of these nine possibilities we compute $n_{?}$ as $2007-S(n)-S(S(n))$ and check whether $S(n_{?})=S(n)$.  
We'll find out that out of the 9 cases, in 4 the value $n_{?}$ has the correct sum of digits.  
This happens for $n_{?}\in \{ 1977, 1980, 1983, 2001 \}$.

 



\textbf{Solution 5}

Claim. The only positive integers $n$ that satisfy the condition are perfect multiples of $3$.

 Proof of claim:We examine the positive integers mod $9$. Here are the cases.

 Case 1. $n \equiv 1 \pmod 9$. Now, we examine $S(n)$ modulo $9$.Case 1.1. The tens digit of $n$ is different from the tens digit of the largest multiple of $9$ under $n$. (In other words, this means we will carry when adding from the perfect multiple of $9$ under $n$.)Observe that when we carry, i.e. Add $1$ onto $1989$ to obtain $1990$, the units digit decreases by $9$ while the tens digit increases by $1$. This means that the sum of the digits decreases by $8$ in total, and we have $-8 \equiv 1 \pmod 9$, so the "mod 9" of the sum increases by $1$. This means that, regardless of whether the sum carries or not, the modulo 9 of the sum of the digits always increases by $1$.

 Case 1.2. The tens digits are the same, which is trivial since the units digit just increases by $1$ which means that the sum is also equivalent to $1 \pmod 9$.

 This means that $S(n) \equiv 1 \pmod 9$ and similarly letting the next $n=S(n)$, $S(S(n)) \equiv 1 \pmod 9$. Summing these, we have $n+S(n)+S(S(n)) \equiv 3 \pmod 9$. Clearly, no integers of this form will satisfy the condition because $2007$ is a perfect multiple of $9$.

 Case 2. $n \equiv 2 \pmod 9$.

 In this case, we apply exactly the same argument. There is at most one carry, which means that the sum of the digits will always be congruent to $2$ mod $9$. Then we can apply similar arguments to get $S(n) \equiv 2 \pmod 9$ and $S(S(n)) \equiv 2 \pmod 9$, so adding gives $n+S(n)+S(S(n)) \equiv 6 \pmod 9$.

 It is trivial to see that for $n \equiv k \pmod 9$, for $0 \leq k \leq 8$, we must have $n+S(n)+S(S(n)) \equiv 3k \pmod 9$. Only when $k=0, 3, 6$ is $3k$ a multiple of $9$, which means that $n$ must be a multiple of $3$.

 Now, we find the integers. Again, consider two cases: Integers that are direct multiples of $9$ and integers that are multiples of $3$ but not $9$.

 Case 1. $n$ is a multiple of $9$. An integer of the form $\overline{20ab}$ will not work since the least such integer is $2007$ which already exceeds our bounds. Thus, we need only consider the integers of the form $\overline{19ab}$. The valid sums of the digits of $n$ are $18$ and $27$ in this case.

 Case 1.1. The sum of the digits is $18$. This means that $S(n)=18, S(S(n))=9$, so $n=2007-18-9=1980$. Clearly this number satisfies our constraints.

 Case 1.2. The sum of the digits is $27$. This means that $S(n)=27, S(S(n))=9$, ,so $n=2007-27-9=1971$. Since the sum of the digits of $1971$ is not $27$, this does not work.

 This means that there is $1$ integer in this case.

 Case 2. $n$ is a multiple of $3$, not $9$..Case 2.1. Integers of the form $\overline{20ab}$. Then $S(n)=3$ or $S(n)=6$; it is trivial to see that $S(n)=6$ exceeds our bounds, so $S(n)=3$ and $n=2007-6=2001$.

 Case 2.2. Integers of the form $\overline{19ab}$. Then $S(n)=12, 15, 21, 24$ and we consider each case separately.

 Case 2.2.1. Integers with $S(n)=12$. That means $n=2007-12-3=1992$ which clearly does not work.

 Case 2.2.2. Integers with $S(n)=15$. That means $n=2007-15-6=1986$ which also does not work

 Case 2.2.3. Integers with $S(n)=21$. That means $n=2007-21-3=1983$ which is valid.

 Case 2.2.4. Integers with $S(n)=24$. That means $n=2007-24-6=1977$ which is also valid.

 We have considered every case, so there are $\boxed{4}$ integers that satisfy the given condition.

 



\textbf{Solution 6 (Rigorous)}

Let the number of digits of $n$ be $m$. If $m = 5$, $n$ will already be greater than $2007$. Notice that $S(n)$ is always at most $9m$. Then if $m = 3$, $n$ will be at most $999$, $S(n)$ will be at most $27$, and $S(S(n))$ will be even smaller than $27$. Clearly we cannot reach a sum of $2007$, unless $m = 4$ (i.e. $n$ has $4$ digits).

 Then, let $n$ be a four digit number in the form $1000a + 100b + 10c + d$. Then $S(n) = a + b + c + d$.

 $S(S(n))$ is the sum of the digits of $a + b + c + d$. We can represent $S(S(n))$ as the sum of the tens digit and the ones digit of $S(n)$. The tens digit in the form of a decimal is

 
$\frac{a + b + c + d}{10}$.

 
To remove the decimal portion, we can simply take the floor of the expression,

 
$\lfloor\frac{a + b + c + d}{10}\rfloor$.

 
Now that we have expressed the tens digit, we can express the ones digit as $S(n) -10$ times the above expression, or

 
$a + b + c + d - 10\lfloor\frac{a + b + c + d}{10}\rfloor$.

 
Adding the two expressions yields the value of $S(S(n))$

 
$= a + b + c + d - 9\lfloor\frac{a + b + c + d}{10}\rfloor$.

 
Combining this expression to the ones for $n$ and $S(n)$ yields

 
$1002a + 102b + 12c + 3d - 9\lfloor\frac{a + b + c + d}{10}\rfloor$.

 
Setting this equal to $2007$ and rearranging a bit yields

 
$12c + 3d = 2007 - 1002a - 102b + 9\lfloor\frac{a + b + c + d}{10}\rfloor$

 $\Rightarrow$ $4c + d = 669 - 334a - 34b + 3\lfloor\frac{a + b + c + d}{10}\rfloor$.

 
(The reason for this slightly weird arrangement will soon become evident)

 
Now we examine the possible values of $a$. If $a \ge 3$, $n$ is already too large. $a$ must also be greater than $0$, or $n$ would be a $3$-digit number. Therefore, $a = 1 \, \text{or} \, 2$. Now we examine by case.

 If $a = 2$, then $b$ and $c$ must both be $0$ (otherwise $n$ would already be greater than $2007$). Substituting these values into the equation yields

 
$d = 1 + 3\lfloor\frac{2 + d}{10}\rfloor$

 $\Rightarrow$ $d=1$.

 
Sure enough, $2001 + (2+1) + 3=2007$.

 Now we move onto the case where $a = 1$. Then our initial equation simplifies to

 
$4c + d = 335 - 34b + 3\lfloor\frac{1 + b + c + d}{10}\rfloor$

 
Since $c$ and $d$ can each be at most $9$, we substitute that value to find the lower bound of $b$. Doing so yields

 
$34b \ge 290 + 3\lfloor\frac{19 + b}{10}\rfloor$.

 
The floor expression is at least $3\lfloor\frac{19}{10}\rfloor=3$ , so the right-hand side is at least $293$. Solving for $b$, we see that $b \ge 9$ $\Rightarrow$ $b=9$. Again, we substitute for $b$ and the equation becomes

 
$4c + d = 29 + 3\lfloor\frac{10 + c + d}{10}\rfloor$

 $\Rightarrow$ $4c + d = 32 + 3\lfloor\frac{c + d}{10}\rfloor$.

 
Just like we did for $b$, we can find the lower bound of $c$ by assuming $d = 9$ and solving:

 
$4c + 9 \ge 29 + 3\lfloor\frac{c + 9}{10}\rfloor$

 $\Rightarrow$ $4c \ge 20 + 3\lfloor\frac{c + 9}{10}\rfloor$

 
The right hand side is $20$ for $c=0$ and $23$ for $c \ge 1$. Solving for c yields $c \ge 6$. Looking back at the previous equation, the floor expression is $0$ for $c+d \le 9$ and $3$ for $c+d \ge 10$. Thus, the right-hand side is $32$ for $c+d \le 9$ and $35$ for $c+d \ge 10$. We can solve these two scenarios as systems of equations/inequalities:

 $4c+d = 32$

 $c+d \le 9$

 and

 $4c+d=35$

 $c+d \ge 10$

 Solving yields three pairs $(c, d):$ $(8, 0)$; $(8, 3)$; and $(7, 7)$. Checking the numbers $1980$, $1983$, and $1977$; we find that all three work. Therefore there are a total of $4$ possibilities for $n$ $\Rightarrow$ $\boxed{\text{D}}$.

 Note: Although this solution takes a while to read (as well as to write) the actual time it takes to think through the process above is very short in comparison to the solution length.

 


\end{enumerate}